\section{Why the Book Needs a Research-Disciplined Case Library}

Without cases, the framework remains too close to a persuasive essay. Without
evidence discipline, the cases risk becoming confident storytelling. This
chapter therefore occupies a deliberately intermediate layer. It does
\emph{not} claim that ordinary study failure, exercise avoidance, or bedtime
drift have already been measured with one-to-one neuroscientific precision.
What it does claim is narrower and more useful:
\begin{enumerate}[nosep]
  \item the variables of the framework provide a stable analysis language;
  \item the research spine from the empirical, emergence, and advanced
        dynamical chapters supplies disciplined mechanism analogies;
  \item application-layer interventions become stronger when they are tied to
        explicit state-transition structure rather than vague motivational talk.
\end{enumerate}

The chapter therefore uses a three-lane reading rule for every case:
\begin{itemize}[nosep]
  \item \textbf{Framework lane:} What are the dominant variables?
  \item \textbf{Mechanism lane:} What research-backed state, path dependence,
        metastability, or attractor ideas are relevant?
  \item \textbf{Evidence lane:} Which claims are strong, which are medium, and
        which would become speculative if pushed too far?
\end{itemize}
